problem:  
Let \( (M^n, g_0) \) be a compact, simply connected Riemannian manifold (\(n \ge 4\)) with **positive sectional curvature**, and let \(g(t)\) denote the **normalized Ricci flow** starting from \(g_0\):  
\[
\frac{\partial g}{\partial t} = -2\,\mathrm{Ric}(g) + \frac{2}{n}r(g)\,g.
\]  
Perelman’s \(\nu\)-functional, defined on the space of Riemannian metrics by  
\[
\nu(g) = \inf_{f,\tau>0} W(g,f,\tau),
\]  
is monotone nondecreasing along the Ricci flow and constant precisely on **shrinking gradient Ricci solitons** (particularly, on Einstein metrics with positive scalar curvature).  

**Open problem:**  
Suppose \(M^n\) admits at least two **non-isometric positive Einstein metrics**, \(g_1\) and \(g_2\), with corresponding Perelman entropies \(\nu(g_1)\) and \(\nu(g_2)\).  

1. Must \(\nu(g_1) \neq \nu(g_2)\)?  
2. If equality \(\nu(g_1) = \nu(g_2)\) were possible, could there exist a **nontrivial ancient solution** \(g(t)\) of the normalized Ricci flow (\(t\in(-\infty,T)\)) satisfying
   \[
   \lim_{t\to -\infty} g(t)=g_1, \quad \lim_{t\to T^-} g(t)=g_2,
   \]
   thereby forming an entropy‑neutral connecting orbit between distinct Einstein metrics?  
3. More broadly, can the space of positive Einstein metrics on a fixed manifold carry multiple distinct **critical points of equal maximal entropy**, implying a flat direction in the entropy landscape that allows analytic connecting trajectories?  

This asks whether Perelman’s entropy functional defines a **strict partial order** among distinct positive Einstein geometries on the same underlying manifold—or whether the entropy landscape can have multiple critical points with identical height.

Why is it a "good" problem:  
This formulation respects the established **gradient‑like structure** of Ricci flow but probes its **fine global geometry** via Perelman’s entropy. The question pushes beyond known monotonicity to ask whether **degeneracies** (equal values for distinct Einstein metrics) can occur, which would profoundly affect how the Ricci flow organizes the space of metrics.  

It directly connects:  
- **Geometric analysis**, through the dynamics of Ricci flow and entropy monotonicity;  
- **Global differential geometry**, through the moduli space of Einstein metrics on positively curved manifolds;  
- **Variational analysis**, through the Morse–theoretic structure of \(\nu\).  

A positive answer would force a reevaluation of entropy rigidity and suggest new analytic families of ancient solutions mediating between Einstein metrics. A negative answer would establish that Perelman’s entropy provides a **strict Morse ordering** on the space of positive Einstein geometries, reinforcing the view of Ricci flow as genuinely gradient‑like.  

The question is simple to state yet touches the deepest structural aspects of Ricci flow and the landscape of Einstein geometries, thereby meeting all criteria for a good and foundational open problem.